
% !TeX root = ../main.tex

\section{Die Schnorr-Identifikation}\label{sec:protokoll}

Das Schnorr-Identifikationsprotokoll ist ein interaktives
kryptografisches Verfahren, mit dem eine Partei die Kenntnis
eines geheimen Schlüssels nachweisen kann, ohne diesen
gegenüber einer anderen Partei offenzulegen.

Das Verfahren basiert auf dem in
Abschnitt~\ref{sec:grundlagen} beschriebenen diskreten
Logarithmusproblem. Es handelt sich um ein sogenanntes
Challenge-Response-Verfahren, bei dem die Authentifizierung
durch den Austausch von drei Nachrichten erfolgt.

Die beweisende Partei wird im Folgenden als \emph{Prover}
und die prüfende Partei als \emph{Verifier} bezeichnet.


\subsection{Parameter und Schlüsselerzeugung}

Die Grundlage des Schnorr-Identifikationsprotokolls bildet
eine endliche zyklische Gruppe
\[
    \G=\langle g\rangle
\]
mit primer Ordnung $q$ und einem Generator $g$.

Für die Realisierung in einer multiplikativen Untergruppe
von $\Z_p^*$ werden zwei Primzahlen $p$ und $q$ gewählt,
für die
\[
    q\mid(p-1)
\]
gilt. Der Generator $g$ besitzt dabei die Ordnung $q$,
sodass
\[
    g^q\equiv 1\pmod p
    \qquad\text{und}\qquad
    g\neq 1
\]
gelten.

Die Parameter $p$, $q$ und $g$ sind öffentlich bekannt
und müssen nicht geheim gehalten werden. Sie sind jedoch
so zu wählen, dass das diskrete Logarithmusproblem in
der verwendeten Gruppe praktisch nicht effizient lösbar
ist.

\paragraph{Schlüsselerzeugung.}

Der Prover erzeugt seinen geheimen Schlüssel $x$, indem
er einen Wert gleichverteilt aus den von null verschiedenen
Elementen von $\Z_q$ auswählt:
\[
    x\xleftarrow{\$}\Z_q\setminus\{0\}.
\]

Anschließend berechnet er den zugehörigen öffentlichen
Schlüssel
\[
    y=g^x.
\]

Bei der konkreten Realisierung in einer Untergruppe von
$\Z_p^*$ erfolgt die Berechnung durch
\[
    y=g^x\bmod p.
\]

Damit besteht das Schlüsselpaar aus
\begin{align*}
    x &\quad\text{(geheimer Schlüssel)},\\
    y &\quad\text{(öffentlicher Schlüssel)}.
\end{align*}

Der geheime Schlüssel verbleibt ausschließlich beim Prover.
Der öffentliche Schlüssel wird dem Verifier zur Verfügung
gestellt und kann beispielsweise während einer Registrierung
einer bestimmten Benutzeridentität zugeordnet werden.

Diese Zuordnung muss vertrauenswürdig erfolgen. Das
Schnorr-Identifikationsprotokoll weist zunächst lediglich
die Kenntnis des zu einem öffentlichen Schlüssel gehörenden
diskreten Logarithmus nach. Eine Aussage über die tatsächliche
Identität einer Person ist nur möglich, wenn der öffentliche
Schlüssel zuvor zuverlässig mit dieser Identität verknüpft
wurde.

Die öffentlichen Parameter und der öffentliche Schlüssel
müssen vor ihrer Verwendung auf Gültigkeit überprüft werden.
Insbesondere muss sichergestellt werden, dass $y$ ein
zulässiges Element der vorgesehenen Untergruppe ist.


\subsection{Ablauf des Drei-Schritt-Protokolls}

Das Schnorr-Identifikationsprotokoll besteht aus drei
aufeinanderfolgenden Nachrichten:

\begin{enumerate}
    \item Commitment des Provers,
    \item Challenge des Verifiers,
    \item Response des Provers.
\end{enumerate}

Ein Protokoll mit dieser Struktur und den später beschriebenen
Sicherheitseigenschaften wird als Sigma-Protokoll
($\Sigma$-Protokoll) bezeichnet.

Zu Beginn der Protokollausführung kennen beide Parteien
die öffentlichen Parameter und den öffentlichen Schlüssel
$y$. Der geheime Schlüssel $x$ ist ausschließlich dem
Prover bekannt.

\subsubsection*{Schritt 1: Commitment}

Der Prover wählt zunächst einen frischen Zufallswert
$r$ gleichverteilt aus $\Z_q$:
\[
    r\xleftarrow{\$}\Z_q.
\]

Dieser temporäre Zufallswert wird auch als Nonce bezeichnet.
Anschließend berechnet der Prover das Commitment
\[
    t=g^r.
\]

Für die Realisierung in $\Z_p^*$ gilt entsprechend
\[
    t=g^r\bmod p.
\]

Der Prover sendet ausschließlich das Commitment $t$ an
den Verifier. Der zugrunde liegende Zufallswert $r$
wird nicht übertragen und muss geheim bleiben.

Das Commitment legt den Prover auf einen bestimmten
Zufallswert fest, bevor ihm die Challenge des Verifiers
bekannt ist.

\subsubsection*{Schritt 2: Challenge}

Nach Erhalt des Commitments überprüft der Verifier dessen
Gültigkeit und wählt anschließend eine Challenge $c$
gleichverteilt aus $\Z_q$:
\[
    c\xleftarrow{\$}\Z_q.
\]

Die Challenge wird an den Prover übertragen.

Entscheidend ist, dass die Challenge erst nach Erhalt
des Commitments gewählt wird und für den Prover im Voraus
nicht vorhersagbar ist.

Dadurch soll verhindert werden, dass ein Teilnehmer ohne
Kenntnis des geheimen Schlüssels bereits vor Beginn
der Prüfung ein passendes Transkript vorbereitet.

\subsubsection*{Schritt 3: Response}

Nach Erhalt der Challenge berechnet der Prover unter
Verwendung seines geheimen Schlüssels $x$ und des zuvor
gewählten Zufallswertes $r$ die Response
\[
    s=(r+cx)\bmod q.
\]

Anschließend sendet der Prover den Wert $s$ an den Verifier.

Der geheime Schlüssel $x$ sowie der temporäre Zufallswert
$r$ werden während des gesamten Protokollablaufs nicht
übertragen.

Die drei Nachrichten lassen sich somit durch folgende
Gleichungen zusammenfassen:
\begin{align}
    t &= g^r
       && \text{Commitment},\\
    c &\xleftarrow{\$}\Z_q
       && \text{Challenge},\\
    s &= (r+cx)\bmod q
       && \text{Response}.
\end{align}

\subsubsection*{Verifikation}

Nach Erhalt der Response besitzt der Verifier sämtliche
für die Prüfung benötigten öffentlichen Werte.

Er überprüft, ob die erhaltenen Werte gültig sind und
ob folgende Gleichung erfüllt ist:
\begin{equation}\label{eq:pruefung}
    g^s\stackrel{?}{=}t\,y^c.
\end{equation}

Im Fall einer multiplikativen Untergruppe von $\Z_p^*$
lautet die Prüfgleichung
\[
    g^s\stackrel{?}{\equiv}t\,y^c\pmod p.
\]

Sind die empfangenen Werte gültig und ist die Gleichung
erfüllt, akzeptiert der Verifier den Nachweis. Andernfalls
wird die Authentifizierung abgelehnt.

\subsubsection*{Darstellung des Nachrichtenflusses}

Der vollständige Nachrichtenfluss zwischen Prover und
Verifier ist in Abbildung~\ref{fig:schnorr-ablauf}
dargestellt.

\begin{figure}[htbp]
    \centering
    \renewcommand{\arraystretch}{1.5}
    \begin{tabular}{ccc}
        \toprule
        \textbf{Prover} &
        &
        \textbf{Verifier}\\
        \midrule
        Kenntnis von $x$ &
        &
        Kenntnis von $y=g^x$\\
        $r\xleftarrow{\$}\Z_q$ &
        & \\
        $t=g^r$ &
        $\xrightarrow{\quad t\quad}$ &
        \\
        &
        $\xleftarrow{\quad c\quad}$ &
        $c\xleftarrow{\$}\Z_q$\\
        $s=(r+cx)\bmod q$ &
        $\xrightarrow{\quad s\quad}$ &
        \\
        &
        &
        $g^s\stackrel{?}{=}t\,y^c$\\
        \bottomrule
    \end{tabular}
    \caption{Nachrichtenfluss des
    Schnorr-Identifikationsprotokolls}
    \label{fig:schnorr-ablauf}
\end{figure}

Die während einer Protokollausführung übertragenen
Nachrichten bilden gemeinsam das sogenannte Transkript
\[
    (t,c,s).
\]

Ein Transkript wird als akzeptierend bezeichnet, wenn
seine Werte gültig sind und die
Prüfgleichung~\eqref{eq:pruefung} erfüllt ist.


\subsection{Korrektheit}

Die Korrektheit des Schnorr-Identifikationsprotokolls
ergibt sich aus den algebraischen Eigenschaften der
verwendeten zyklischen Gruppe.

Es wird angenommen, dass beide Parteien das Protokoll
ordnungsgemäß ausführen und der Prover den geheimen
Schlüssel $x$ kennt.

Aus der Definition des öffentlichen Schlüssels
\[
    y=g^x
\]
sowie des Commitments
\[
    t=g^r
\]
und der Response
\[
    s=(r+cx)\bmod q
\]
folgt aufgrund der Gruppenordnung $q$:
\begin{align}
    g^s
        &=g^{r+cx}\\
        &=g^r\cdot g^{cx}\\
        &=g^r\cdot(g^x)^c\\
        &=t\cdot y^c.
\end{align}

Somit ist die Prüfgleichung
\[
    g^s=t\cdot y^c
\]
für einen ehrlich erzeugten Protokollablauf erfüllt.

Die Reduktion der Response modulo $q$ verändert das
resultierende Gruppenelement nicht, da $g^q=e$ gilt.

Damit ist gezeigt, dass ein Prover, der den geheimen
Schlüssel kennt und das Protokoll korrekt ausführt,
die Prüfung eines ehrlichen Verifiers erfolgreich
bestehen kann.

Diese Eigenschaft wird als Vollständigkeit
(\emph{Completeness}) bezeichnet.


\subsection{Numerisches Beispiel}

Zur Veranschaulichung des Protokollablaufs werden die
bereits in Abschnitt~\ref{sec:grundlagen} eingeführten
Parameter verwendet:
\[
    p=23,\qquad q=11,\qquad g=2.
\]

Die Gruppe $\G=\langle 2\rangle$ besitzt modulo $23$
die Ordnung $11$.

\paragraph{Schlüsselerzeugung.}

Der Prover wählt den geheimen Schlüssel
\[
    x=7.
\]

Der zugehörige öffentliche Schlüssel ergibt sich aus
\begin{align*}
    y&=g^x\bmod p\\
     &=2^7\bmod 23\\
     &=13.
\end{align*}

\paragraph{Commitment.}

Für die Protokollausführung wird beispielhaft der
Zufallswert
\[
    r=4
\]
gewählt.

Das Commitment lautet damit
\begin{align*}
    t&=g^r\bmod p\\
     &=2^4\bmod 23\\
     &=16.
\end{align*}

Der Prover sendet den Wert $t=16$ an den Verifier.

\paragraph{Challenge.}

Der Verifier wählt beispielhaft die Challenge
\[
    c=3
\]
und sendet diese an den Prover.

\paragraph{Response.}

Der Prover berechnet die Response
\begin{align*}
    s&=(r+cx)\bmod q\\
     &=(4+3\cdot 7)\bmod 11\\
     &=25\bmod 11\\
     &=3.
\end{align*}

Der berechnete Wert $s=3$ wird an den Verifier übertragen.

\paragraph{Verifikation.}

Der Verifier überprüft die Gleichung
\[
    g^s\stackrel{?}{\equiv}t\,y^c\pmod p.
\]

Für die linke Seite ergibt sich
\[
    g^s=2^3\bmod 23=8.
\]

Die rechte Seite liefert
\begin{align*}
    t\,y^c\bmod p
        &=16\cdot 13^3\bmod 23\\
        &=35152\bmod 23\\
        &=8.
\end{align*}

Da beide Seiten übereinstimmen, akzeptiert der Verifier
den Nachweis:
\[
    8=8.
\]

Die während der Protokollausführung verwendeten Werte
sind in Tabelle~\ref{tab:schnorr-beispiel} zusammengefasst.

\begin{table}[htbp]
    \centering
    \begin{tabular}{lrl}
        \toprule
        Symbol & Wert & Bedeutung\\
        \midrule
        $p$ & 23 & Primzahl als Modulus\\
        $q$ & 11 & Gruppenordnung\\
        $g$ & 2  & Generator\\
        $x$ & 7  & Geheimer Schlüssel\\
        $y$ & 13 & Öffentlicher Schlüssel\\
        $r$ & 4  & Temporärer Zufallswert\\
        $t$ & 16 & Commitment\\
        $c$ & 3  & Challenge\\
        $s$ & 3  & Response\\
        \bottomrule
    \end{tabular}
    \caption{Beispiel einer erfolgreichen
    Schnorr-Identifikation}
    \label{tab:schnorr-beispiel}
\end{table}

Das resultierende akzeptierende Transkript lautet
\[
    (t,c,s)=(16,3,3).
\]

Der geheime Schlüssel $x$ wurde während der
Protokollausführung nicht direkt an den Verifier übertragen.

Aufgrund der bewusst kleinen Parameter könnte der geheime
Schlüssel in diesem Beispiel allerdings durch Berechnung
des diskreten Logarithmus leicht bestimmt werden.
Das Beispiel dient daher ausschließlich der Veranschaulichung
und stellt keine kryptografisch sichere Implementierung dar.


\subsection{Sicherheitsidee und Grenzen}

Das Schnorr-Identifikationsprotokoll besitzt drei
wesentliche Eigenschaften, die für Sigma-Protokolle
charakteristisch sind:

\begin{itemize}
    \item Vollständigkeit (\emph{Completeness}),
    \item spezielle Wissensextraktion
          (\emph{Special Soundness}),
    \item Zero Knowledge gegenüber einem ehrlichen
          Verifier (\emph{Honest-Verifier Zero Knowledge}).
\end{itemize}

Diese Eigenschaften beschreiben unterschiedliche Aspekte
der Protokollsicherheit und müssen voneinander
unterschieden werden.

\subsubsection*{Vollständigkeit}

Die Vollständigkeit besagt, dass ein ehrlicher Prover,
der den geheimen Schlüssel kennt, einen ehrlichen
Verifier von dieser Kenntnis überzeugen kann.

Wie im Abschnitt zur Korrektheit gezeigt wurde, gilt
für eine ordnungsgemäße Protokollausführung stets
\[
    g^s=t\,y^c.
\]

Unter der Voraussetzung gültiger Parameter und einer
korrekten Ausführung akzeptiert der Verifier daher
den Nachweis mit Wahrscheinlichkeit $1$.

Die Vollständigkeit ist somit eine Aussage über das
Verhalten ehrlicher Teilnehmer und noch kein Nachweis
dafür, dass sich ein Angreifer nicht erfolgreich
authentifizieren kann.


\subsubsection*{Spezielle Wissensextraktion}

Die spezielle Wissensextraktion beschreibt eine
Eigenschaft, die für den Nachweis der Kenntnis des
geheimen Schlüssels wesentlich ist.

Angenommen, es liegen zwei akzeptierende Transkripte
\[
    (t,c_1,s_1)
    \qquad\text{und}\qquad
    (t,c_2,s_2)
\]
mit demselben Commitment $t$, jedoch unterschiedlichen
Challenges
\[
    c_1\neq c_2
\]
vor.

Da beide Transkripte akzeptierend sind, gilt
\begin{align}
    g^{s_1}&=t\,y^{c_1},\\
    g^{s_2}&=t\,y^{c_2}.
\end{align}

Durch Multiplikation der ersten Gleichung mit dem
Inversen der zweiten Gleichung ergibt sich
\[
    g^{s_1-s_2}=y^{c_1-c_2}.
\]

Mit $y=g^x$ folgt
\[
    g^{s_1-s_2}=g^{x(c_1-c_2)}.
\]

Da $g$ die Ordnung $q$ besitzt, müssen die Exponenten
modulo $q$ übereinstimmen:
\[
    s_1-s_2\equiv x(c_1-c_2)\pmod q.
\]

Weil $q$ prim ist und $c_1\neq c_2$ gilt, besitzt
$c_1-c_2$ ein multiplikatives Inverses modulo $q$.

Durch Multiplikation mit diesem Inversen erhält man
\begin{equation}\label{eq:schnorr-extraktion}
    x=(s_1-s_2)(c_1-c_2)^{-1}\bmod q.
\end{equation}

Aus zwei akzeptierenden Transkripten mit identischem
Commitment und unterschiedlichen Challenges kann somit
der geheime Schlüssel effizient berechnet werden.

Diese Eigenschaft wird in einem theoretischen
Sicherheitsbeweis durch einen sogenannten
Wissensextraktor genutzt.

Der Wissensextraktor betrachtet eine Protokollausführung
und versetzt den untersuchten Prover gedanklich in den
Zustand unmittelbar nach dem Senden des Commitments
zurück. Anschließend wird eine andere Challenge
verwendet.

Kann der Prover mit hinreichend hoher Wahrscheinlichkeit
gültige Antworten auf unterschiedliche Challenges
für dasselbe Commitment erzeugen, lässt sich daraus
der zum öffentlichen Schlüssel gehörende diskrete
Logarithmus extrahieren.

Dieses Verfahren wird als \emph{Rewinding} bezeichnet.
Es handelt sich um eine Technik des Sicherheitsbeweises
und nicht um einen regulären Schritt des eigentlichen
Identifikationsprotokolls.

Die spezielle Wissensextraktion bildet damit eine
Grundlage für die Interpretation des Schnorr-Protokolls
als Wissensnachweis.

Für eine gleichverteilt aus $\Z_q$ gewählte Challenge
beträgt der sogenannte Wissensfehler des Sigma-Protokolls
\[
    \varepsilon=\frac{1}{q}.
\]

Ein Teilnehmer kann beispielsweise im Voraus eine
Challenge erraten und ein passendes Transkript
konstruieren. Unter der Annahme einer unabhängigen
und gleichverteilten Challenge beträgt die
Erfolgswahrscheinlichkeit einer solchen Strategie
genau $1/q$.

Ein Prover, dessen Erfolgswahrscheinlichkeit wesentlich
über diesem Wissensfehler liegt, kann unter den
Voraussetzungen des entsprechenden Extraktionsarguments
zur Berechnung eines gültigen geheimen Schlüssels
verwendet werden.

Bei ausreichend großer Gruppenordnung und unter der
Annahme der Schwierigkeit des diskreten Logarithmusproblems
ergibt sich daraus die Sicherheitsidee des
Identifikationsverfahrens.


\subsubsection*{Honest-Verifier Zero Knowledge}

Die Zero-Knowledge-Eigenschaft beschreibt, dass ein
Nachweis keine zusätzlichen Informationen über das
geheime Wissen des Provers offenlegen soll.

Für das Schnorr-Identifikationsprotokoll lässt sich
diese Eigenschaft zunächst für einen ehrlichen Verifier
zeigen.

Ein ehrlicher Verifier hält sich an die
Protokollspezifikation und wählt seine Challenge
insbesondere unabhängig und gleichverteilt aus
dem vorgesehenen Wertebereich.

Die Beweisidee besteht darin, einen Simulator zu
konstruieren, der ein akzeptierendes Transkript
erzeugen kann, ohne den geheimen Schlüssel $x$
zu kennen.

Der Simulator kennt lediglich die öffentlichen Parameter
und den öffentlichen Schlüssel $y$.

Er führt folgende Schritte aus:
\begin{enumerate}
    \item Wähle eine Challenge $c$ gleichverteilt
    aus $\Z_q$:
    \[
        c\xleftarrow{\$}\Z_q.
    \]

    \item Wähle eine Response $s$ unabhängig und
    gleichverteilt aus $\Z_q$:
    \[
        s\xleftarrow{\$}\Z_q.
    \]

    \item Berechne das Commitment rückwärts aus
    der Prüfgleichung:
    \[
        t=g^s\cdot y^{-c}.
    \]
\end{enumerate}

Der Ausdruck $y^{-c}$ bezeichnet dabei die entsprechende
inverse Potenz des Gruppenelements $y$.

Der Simulator gibt anschließend das Transkript
\[
    (t,c,s)
\]
aus.

Dieses Transkript erfüllt die Prüfgleichung, denn
\begin{align*}
    t\,y^c
        &=(g^s\cdot y^{-c})\cdot y^c\\
        &=g^s.
\end{align*}

Es handelt sich somit um ein akzeptierendes Transkript,
obwohl der Simulator den geheimen Schlüssel nicht
kennt.

Darüber hinaus besitzen die simulierten Transkripte
dieselbe Wahrscheinlichkeitsverteilung wie die
Transkripte einer ehrlichen Protokollausführung.

In einer regulären Protokollausführung gilt für jede
fest gewählte Challenge $c$
\[
    s=(r+cx)\bmod q.
\]

Da $r$ gleichverteilt aus $\Z_q$ gewählt wird, ist
auch $s$ gleichverteilt in $\Z_q$. Die Addition des
festen Wertes $cx$ verändert diese Gleichverteilung
nicht.

Das zugehörige Commitment lässt sich durch
\[
    t=g^r
     =g^{s-cx}
     =g^s\cdot y^{-c}
\]
darstellen.

Damit stimmen die Verteilungen der simulierten und
der ehrlich erzeugten Transkripte exakt überein.

Dies wird als perfekte Honest-Verifier-Zero-Knowledge-
Eigenschaft bezeichnet.

Der Verifier erhält aus einer ehrlichen
Protokollausführung folglich keine zusätzliche
Information über den geheimen Schlüssel, die nicht
auch aus den öffentlichen Informationen simuliert
werden könnte.

Dabei ist zu beachten, dass der öffentliche Schlüssel
selbst bereits den geheimen Schlüssel mathematisch
festlegt. Die praktische Geheimhaltung des Schlüssels
beruht weiterhin auf der Schwierigkeit des diskreten
Logarithmusproblems.

\paragraph{Abgrenzung zur Wissenseigenschaft.}

Die Möglichkeit, ein akzeptierendes Transkript ohne
Kenntnis des geheimen Schlüssels zu simulieren, steht
nicht im Widerspruch zur Wissenseigenschaft des
Protokolls.

Der Simulator darf zunächst die Challenge und die
Response auswählen und erst anschließend das passende
Commitment berechnen.

Während einer tatsächlichen Protokollausführung muss
der Prover das Commitment jedoch bereits übertragen,
bevor die Challenge des Verifiers bekannt ist.

Die zeitliche Reihenfolge der Nachrichten ist daher
ein wesentlicher Bestandteil der Protokollsicherheit.

Ein einzelnes akzeptierendes Transkript stellt deshalb
für sich allein keinen unabhängig überprüfbaren Beweis
dar, dass eine bestimmte Person tatsächlich den geheimen
Schlüssel kennt.


\subsubsection*{Bedeutung frischer Zufallswerte}

Die Sicherheit des Schnorr-Identifikationsprotokolls
setzt voraus, dass für jede Protokollausführung ein
neuer, unabhängiger und kryptografisch sicherer
Zufallswert $r$ erzeugt wird.

Wird derselbe Zufallswert in zwei Ausführungen verwendet,
so entsteht dasselbe Commitment
\[
    t=g^r.
\]

Erhält der Prover dabei zwei unterschiedliche Challenges
$c_1$ und $c_2$, ergeben sich die Responses
\begin{align}
    s_1&=(r+c_1x)\bmod q,\\
    s_2&=(r+c_2x)\bmod q.
\end{align}

Durch Subtraktion der beiden Gleichungen fällt der
Zufallswert $r$ weg:
\[
    s_1-s_2\equiv(c_1-c_2)x\pmod q.
\]

Der geheime Schlüssel kann daher unmittelbar mithilfe
von Gleichung~\eqref{eq:schnorr-extraktion} berechnet
werden.

Die Wiederverwendung eines Zufallswertes kann somit
zur vollständigen Kompromittierung des langfristigen
geheimen Schlüssels führen.

\paragraph{Numerisches Beispiel einer
Nonce-Wiederverwendung.}

Es werden erneut die Beispielparameter
\[
    p=23,\qquad q=11,\qquad g=2,\qquad x=7
\]
verwendet.

Angenommen, der Prover verwendet in zwei verschiedenen
Protokollausführungen denselben Zufallswert
\[
    r=4.
\]

Das Commitment lautet in beiden Ausführungen
\[
    t=2^4\bmod 23=16.
\]

Für die erste Challenge $c_1=3$ erhält man
\begin{align*}
    s_1&=(4+3\cdot 7)\bmod 11\\
       &=3.
\end{align*}

Für die zweite Challenge $c_2=5$ ergibt sich
\begin{align*}
    s_2&=(4+5\cdot 7)\bmod 11\\
       &=6.
\end{align*}

Damit liegen die beiden akzeptierenden Transkripte
\[
    (16,3,3)
    \qquad\text{und}\qquad
    (16,5,6)
\]
vor.

Der geheime Schlüssel kann nun durch
\begin{align*}
    x
        &=(s_1-s_2)(c_1-c_2)^{-1}\bmod q\\
        &=(3-6)(3-5)^{-1}\bmod 11\\
        &=(-3)(-2)^{-1}\bmod 11
\end{align*}
berechnet werden.

Da
\[
    (-2)^{-1}\equiv 5\pmod{11}
\]
gilt, ergibt sich
\begin{align*}
    x&=(-3)\cdot 5\bmod 11\\
     &=-15\bmod 11\\
     &=7.
\end{align*}

Der ursprüngliche geheime Schlüssel wurde damit
vollständig rekonstruiert.

Dieses Beispiel verdeutlicht, dass die Sicherheit
des Protokolls nicht ausschließlich von der Schwierigkeit
des diskreten Logarithmusproblems abhängt, sondern
auch von der Qualität und korrekten Verwendung der
erzeugten Zufallswerte.

Auch die Offenlegung oder Vorhersagbarkeit eines
einzelnen Zufallswertes kann den geheimen Schlüssel
gefährden. Für eine von null verschiedene Challenge
$c$ gilt bei bekanntem $r$
\[
    x=(s-r)c^{-1}\bmod q.
\]

Der temporäre Zufallswert muss daher sowohl während
als auch nach seiner Verwendung vor unbefugtem
Zugriff geschützt werden. Insbesondere darf ein
bereits verwendetes Commitment nicht für die
Beantwortung weiterer Challenges wiederverwendet
werden.


\subsubsection*{Angreifermodell und weitere
Sicherheitsgrenzen}

Die beschriebenen mathematischen Eigenschaften gelten
unter bestimmten Voraussetzungen. Sie erlauben keine
pauschale Aussage über die Sicherheit jeder praktischen
Implementierung.

Insbesondere sind folgende Aspekte zu berücksichtigen:

\begin{description}

    \item[Unvorhersehbare Challenges:]
    Die Challenge muss unabhängig vom Prover und
    mithilfe eines kryptografisch sicheren
    Zufallszahlengenerators erzeugt werden.
    Ist die Challenge vorhersehbar, kann ein Angreifer
    ein passendes Commitment vorbereiten, ohne den
    geheimen Schlüssel zu kennen.

    \item[Validierung der Gruppenparameter:]
    Die verwendeten Parameter müssen die notwendigen
    mathematischen Voraussetzungen erfüllen.
    Dazu zählen insbesondere die Primzahleigenschaften
    von $p$ und $q$, die Teilbarkeitsbedingung
    $q\mid(p-1)$ sowie die korrekte Ordnung des
    Generators $g$.

    \item[Validierung der Gruppenelemente:]
    Öffentliche Schlüssel und Commitments müssen
    auf ihre Zugehörigkeit zur vorgesehenen Untergruppe
    überprüft werden.

    Für ein Element $u$ in der multiplikativen
    Untergruppe müssen insbesondere
    \[
        1\leq u\leq p-1
    \]
    und
    \[
        u^q\equiv 1\pmod p
    \]
    gelten.

    Für den öffentlichen Schlüssel ist zusätzlich
    $y\neq 1$ erforderlich, da der zugehörige
    diskrete Logarithmus andernfalls öffentlich
    bekannt wäre.

    \item[Replay-Angriffe:]
    Ein aufgezeichnetes Transkript kann nicht ohne
    Weiteres für eine spätere Authentifizierung
    verwendet werden, da der Verifier eine neue,
    unabhängige Challenge erzeugt.

    Stimmt die neue Challenge mit der aufgezeichneten
    Challenge überein, ist ein Replay jedoch grundsätzlich
    möglich. Bei gleichverteilter Auswahl aus $\Z_q$
    beträgt die Wahrscheinlichkeit hierfür $1/q$.
    Die Challenge muss daher einen ausreichend
    großen Wertebereich besitzen.

    \item[Relay-Angriffe:]
    Eine frische Challenge verhindert nicht zwangsläufig
    die Weiterleitung einer laufenden Authentifizierung.
    Ein Angreifer könnte Nachrichten zwischen einem
    legitimen Prover und einem Verifier weiterleiten,
    ohne den geheimen Schlüssel selbst zu kennen.

    Zusätzliche Mechanismen zur Bindung der
    Authentifizierung an den Kommunikationskanal,
    die Gegenstelle oder den Anwendungskontext
    können daher erforderlich sein.

    \item[Manipulierte Verifier:]
    Die hier gezeigte Zero-Knowledge-Eigenschaft
    bezieht sich auf einen ehrlichen Verifier.
    Daraus folgt nicht unmittelbar eine entsprechende
    Sicherheit gegenüber beliebig vom Protokoll
    abweichenden Verifiern.

    Weitergehende Aussagen erfordern ein präzises
    Sicherheitsmodell und entsprechende Beweise.

    \item[Implementierungssicherheit:]
    Fehler bei der Implementierung, unsichere
    Zufallszahlengeneratoren, Seitenkanalangriffe
    oder ein unzureichender Schutz des geheimen
    Schlüssels können die Sicherheit des Verfahrens
    beeinträchtigen.

    \item[Quantencomputer:]
    Die langfristige Sicherheit des Verfahrens
    hängt von der Schwierigkeit des diskreten
    Logarithmusproblems ab.
    Ein hinreichend leistungsfähiger, fehlertoleranter
    Quantencomputer könnte dieses Problem mithilfe
    des Shor-Algorithmus effizient lösen.

    Das klassische Schnorr-Identifikationsprotokoll
    ist daher kein postquantenkryptografisches
    Verfahren.

\end{description}


\subsection{Abgrenzung zu nichtinteraktiven Verfahren}

Das in diesem Abschnitt beschriebene
Schnorr-Identifikationsprotokoll setzt eine aktive
Interaktion zwischen Prover und Verifier voraus.

Eine bekannte Erweiterung besteht darin, die zufällige
Challenge des Verifiers durch einen aus öffentlichen
Informationen berechneten kryptografischen Hashwert
zu ersetzen. Dieser Ansatz wird als
Fiat-Shamir-Transformation bezeichnet.

Abstrakt kann eine solche Challenge beispielsweise
durch
\[
    c=H(g,y,t,\text{Kontext})\bmod q
\]
dargestellt werden.

Der Prover kann dadurch die Challenge selbst berechnen,
ohne zuvor eine Nachricht des Verifiers erhalten
zu müssen.

Die konkrete Definition und Kodierung der Hash-Eingaben
sowie die Abbildung des Hashwertes in den vorgesehenen
Challenge-Bereich sind dabei sicherheitsrelevant.

Diese Transformation ermöglicht unter geeigneten
Annahmen die Konstruktion nichtinteraktiver
Wissensnachweise und bildet außerdem eine Grundlage
für Schnorr-Signaturverfahren.

Bei einer digitalen Signatur werden zusätzlich die
zu signierende Nachricht und gegebenenfalls weitere
Kontextinformationen in die Berechnung der Challenge
einbezogen.

Die Sicherheit einer solchen Erweiterung muss jedoch
gesondert untersucht werden. Die Eigenschaften des
interaktiven Protokolls können nicht ohne weitere
Annahmen unmittelbar auf jede nichtinteraktive
Variante übertragen werden.


\subsection{Zusammenfassung}

Das Schnorr-Identifikationsprotokoll ermöglicht den
Nachweis der Kenntnis eines geheimen diskreten
Logarithmus, ohne diesen direkt an den Verifier
zu übertragen.

Die Kombination aus Commitment, Challenge und
Response gewährleistet bei einer korrekten
Protokollausführung die Vollständigkeit des Verfahrens.

Die spezielle Wissensextraktion zeigt, dass aus
zwei akzeptierenden Transkripten mit identischem
Commitment und unterschiedlichen Challenges der
geheime Schlüssel berechnet werden kann.

Gleichzeitig lässt sich durch Simulation nachweisen,
dass eine ehrliche Protokollausführung keine
zusätzlichen Informationen über den geheimen
Schlüssel offenlegt.

Die Sicherheit des Verfahrens beruht somit auf dem
Zusammenspiel der mathematischen Eigenschaften des
Protokolls, der Schwierigkeit des diskreten
Logarithmusproblems und der korrekten praktischen
Implementierung.
